\section{Deployment guidelines}
\label{sec:guidelines}

The results point to four practices for operators of synchronous call graphs. Each practice follows from a measured effect and carries the limit that the evidence places on it.

\textbf{First, size and alert on the thread-pool utilization of the entry service, not on host CPU.} On the chains, the root pool reached utilization $0.75$ to $0.79$ while the processes used $9.2\%$ to $10.3\%$ of the eight cores (Table~\ref{tab:knee}). Host CPU was therefore a poor early warning. Pool utilization is computed from the holding time of the entry service and the arrival rate, as in Eq.~\eqref{eq:rho}. For the asynchronous and hybrid runtimes the same quantity is a holding-time measure, and it should not be read as occupancy.

\textbf{Second, treat call-graph depth and in-cast width as capacity.} The saturation horizon of a depth-eight chain fell to about a quarter of the depth-two horizon, from $312$ to $77$\,req/s in our testbed. A width-eight diamond saturated between $60$ and $80$\,req/s, while a width-four diamond of the same depth saturated between $80$ and $120$\,req/s. Admission control belongs at the entry service and should be sized from a measured horizon for the call graph in question. Where a service has replicas, choosing the less loaded of two random replicas removes most of the imbalance at negligible cost~\cite{mitzenmacher2001power}, but that policy was not evaluated here. The horizons are specific to a single eight-core host with 10\,ms local work, and they must be remeasured for any other deployment.

\textbf{Third, match the control to the fault.} Deadline-budget shedding protected successful requests when the fault consumed capacity. Under a one-second hang or a four-fold slowdown at high load, DAC-Shed kept the P99.9 of successful requests at or below $256\,\mathrm{ms}$, where the static breaker and naive retries reached several seconds. Under transient errors the picture reverses, and naive retries delivered the most goodput. A bounded, budget-gated retry is the candidate for transient errors, and the offline replay recovered $69.2\%$ of the first-attempt errors. That replay is not a live result, and it should be tested live before it is adopted.

\textbf{Fourth, budget the sidecar in CPU and latency separately, and bound the tracing cost before enabling it at high load.} In the emulation, the sidecar used $23\%$ to $32\%$ of the mesh CPU and added $1.8$ to $3.4\,\mathrm{ms}$ per hop, while serialization was below one per cent of CPU for these payloads. Tracing at sampling rates of $0.25$ or higher raised median P99 by $4.3$ to $10.6\,\mathrm{ms}$ at high load. That difference is suggestive and confounded with background load, so it should be measured on the production stack before a sampling policy is set. The sidecar figures describe an emulated proxy and must be remeasured on the production proxy before they are used for capacity planning.

Two cautions apply to all four practices. The runtime choice matters, because the synchronous runtime saturated at the thread pool while the asynchronous and hybrid runtimes did not, and the hybrid runtime was markedly slower on the in-cast. The choice of runtime should therefore be made with the call-graph shape in view. Although Section~\ref{sec:gate} demonstrates that saturation horizons and controller rankings remain stable across all phases up to 1.2 cores of background load, operators running latency-critical benchmarks should still isolate measurement hosts from concurrent workloads.
